\section{Introducción: de la reconstrucción 3D a la generación 3D generativa}

Philipp Henzler lleva investigando en el ámbito de la reconstrucción 3D y la generación desde hace tiempo, aún durante su doctorado en 2018. Sus principales trabajos incluyen la reconstrucción 3D autosupervisada (self-supervised 3D reconstruction) y el modelado de apariencia (appearance modeling). Tras unirse a Google en 2022, se centró en la reconstrucción 3D generativa y el relighting 3D. Su equipo logró aplicar los resultados de sus investigaciones al producto Google Shopping: primero con la reconstrucción 3D de calzado y luego, mediante modelos de video, extendiéndolo a cualquier categoría de objetos.

El tema central de esta conferencia es: \textbf{cómo combinar los métodos de generación 3D con los modelos de video generales}, y por qué el modelado 3D explícito sigue siendo indispensable hoy en día, cuando los modelos de video son cada vez más potentes.

\begin{importantbox}{Preguntas clave de la conferencia}
Si los modelos de video en 2D ya pueden generar videos altamente realistas, ¿por qué todavía necesitamos modelar 3D de forma explícita? ¿Cuáles son las ventajas y desventajas respectivas de los métodos 3D y de los modelos generativos en 2D? ¿Cómo se pueden combinar ambos para obtener los mejores resultados?
\end{importantbox}
